\section{The finite arithmetic certificate}
\label{app:certificate}

The following complete Python program verifies the finite calculation in Lemma~\ref{seed:product}. It uses only integer arithmetic and the Python standard library. In the dictionaries representing valuation columns, $(r,p)$ denotes the valuation at the prime $p$ of ratio number $r$. The key $(-1,0)$ represents the appended row of ones. The rank routine performs Gaussian elimination over $\F_{1009}$, storing one normalized pivot for each leading row. Thus its return value is exactly the rank of the augmented matrix. Separately, $1009$ is prime: none of $2,3,5,7,11,13,17,19,23,29,31$ divides it, and $32^2>1009$.

The sieve stores a prime divisor of every integer greater than one. Repeated division therefore computes its prime valuations exactly. The preliminary cutoff and deletion of singly supported columns give the matrix used in the proof; every deleted column receives coefficient zero in the eventual relation. The final assertion is the sole finite calculation needed for existence of the two product identities.

\newpage
\lstinputlisting{../verification/support/supplied_finite_rank.py}

\noindent An explicit alternative relation is supplied in
\nolinkurl{verification/support/explicit-relation.json}. It has $849$ nonzero integer
coefficients, each with at most $140$ decimal digits, and positive
coefficient sum. The standard-library program
\nolinkurl{verification/validation/explicit_relation.py} verifies both product
identities by trial division of all six integers in every row, without
using the matrix or its rank. It also checks every decreasing-remainder
chain and the variable counts in the polynomial construction. This
independent witness is supplementary: the rank computation above and the
adjugate argument already prove Lemma~\ref{seed:product}.
